\documentclass[11pt,a4paper,openany]{report}

% Fichier autonome : téléverser uniquement pnr.tex dans Overleaf.
\usepackage[T1]{fontenc}
\usepackage[utf8]{inputenc}
\usepackage[french]{babel}
\usepackage{lmodern}
\usepackage{microtype}
\usepackage[a4paper,margin=24mm,headheight=15pt]{geometry}
\usepackage{xcolor}
\usepackage{booktabs}
\usepackage{longtable}
\usepackage{enumitem}
\usepackage{amsmath}
\usepackage{listings}
\usepackage[most]{tcolorbox}
\usepackage{fancyhdr}
\usepackage{hyperref}
\usepackage{bookmark}

\definecolor{GuideBlue}{HTML}{174A7E}
\definecolor{GuideGreen}{HTML}{237A57}
\definecolor{GuideOrange}{HTML}{B45F06}
\definecolor{GuideRed}{HTML}{A61B1B}
\definecolor{GuideGray}{HTML}{F4F5F7}

\hypersetup{
  colorlinks=true,
  linkcolor=GuideBlue,
  urlcolor=GuideBlue,
  pdftitle={Implémentation physique avec Innovus},
  pdfauthor={Karim Sabra}
}

\pagestyle{fancy}
\fancyhf{}
\fancyhead[R]{\small\nouppercase{\leftmark}}
\fancyfoot[C]{\thepage}
\renewcommand{\headrulewidth}{0.4pt}
\setcounter{tocdepth}{0}
\setlist{nosep,leftmargin=*}

\lstdefinelanguage{TclGuide}{
  morekeywords={set,proc,if,elseif,else,foreach,return,error,catch,puts,exit,
    source,global,file,info,dict,string,open,close,lappend,init_design,
    globalNetConnect,applyGlobalNets,checkDesign,floorPlan,defOut,placeDesign,
    checkPlace,setDesignMode,routeDesign,saveDesign,saveNetlist,redirect,
    create_constraint_mode,create_library_set,create_rc_corner,
    create_delay_corner,create_analysis_view,set_analysis_view,addRing,sroute,
    verifyConnectivity,create_ccopt_clock_tree_spec,ccopt_design,optDesign,
    verify_drc,verifyProcessAntenna,timeDesign,extractRC,rcOut,
    streamOut,saveNetlist,addFiller,report_timing},
  sensitive=true,
  morecomment=[l]{\#},
  morestring=[b]"
}

\lstdefinestyle{guidecode}{
  basicstyle=\ttfamily\scriptsize,
  keywordstyle=\color{GuideBlue}\bfseries,
  commentstyle=\color{GuideGreen},
  stringstyle=\color{GuideOrange},
  backgroundcolor=\color{GuideGray},
  frame=single,
  rulecolor=\color{black!20},
  columns=fullflexible,
  keepspaces=true,
  showstringspaces=false,
  breaklines=true,
  tabsize=2,
  numbers=left,
  numberstyle=\tiny\color{black!50},
  xleftmargin=1.5em,
  literate=
    {à}{{\`a}}1 {ç}{{\c c}}1 {é}{{\'e}}1 {è}{{\`e}}1
    {ê}{{\^e}}1 {î}{{\^i}}1 {ô}{{\^o}}1 {ù}{{\`u}}1
}

\newtcolorbox{keybox}[1][À retenir]{
  breakable,colback=GuideBlue!5,colframe=GuideBlue,title=\textbf{#1}}
\newtcolorbox{warningbox}[1][Attention]{
  breakable,colback=GuideOrange!7,colframe=GuideOrange,title=\textbf{#1}}
\newtcolorbox{failurebox}[1][Diagnostic]{
  breakable,colback=GuideRed!5,colframe=GuideRed,title=\textbf{#1}}

\begin{document}

\begin{titlepage}
  \centering
  \vspace*{28mm}
  {\Huge\bfseries\color{GuideBlue}Implémentation physique avec Innovus\par}
  \vspace{7mm}
  {\LARGE De la netlist placée au routage vérifié\par}
  \vfill
  {\Large Karim Sabra\par}
  \vspace{4mm}
  {Guide autonome compatible Overleaf\par}
  {\today\par}
\end{titlepage}

\chapter*{Mode d'emploi Overleaf}
\addcontentsline{toc}{chapter}{Mode d'emploi Overleaf}

Ce document ne dépend d'aucun autre fichier LaTeX. Dans Overleaf, créer un
projet vide, téléverser uniquement \texttt{pnr.tex}, le sélectionner comme
document principal et choisir le compilateur pdfLaTeX.

\begin{warningbox}[Portée]
Les noms de couches, sites, pins d'alimentation, cellules et options physiques
dépendent du PDK. Les exemples illustrent la méthode; ils doivent être adaptés
aux fichiers autorisés du laboratoire. Une commande Innovus réussie n'est pas
une preuve de fermeture physique.
\end{warningbox}

\tableofcontents

\chapter{But et chaîne d'implémentation}

\section{Entrées nécessaires}

\begin{longtable}{p{0.23\textwidth}p{0.30\textwidth}p{0.37\textwidth}}
\toprule
Entrée & Exemple & Rôle \\
\midrule
Netlist Genus & \texttt{full\_adder\_comb.mapped.v} & logique mappée \\
Contraintes & \texttt{mapped.sdc} & intention de timing \\
Technology LEF & \texttt{technology.lef} & couches et règles abstraites \\
Standard-cell LEF & \texttt{stdcells.lef} & géométrie et pins des cellules \\
Liberty & BC, TC, WC & arcs de timing et puissance \\
QRC technology & un fichier par corner RC & extraction des interconnexions \\
Stream map & \texttt{stream.map} & correspondance couches vers GDS \\
GDS des cellules & \texttt{stdcells.gds} & géométrie de référence \\
\bottomrule
\end{longtable}

\section{Étapes}

\begin{enumerate}
  \item préflight des entrées et de la cohérence logique/physique;
  \item import de la netlist, des LEF, du MMMC et du SDC;
  \item floorplan, rows, pins et réseau d'alimentation;
  \item placement et optimisation pré-CTS;
  \item CTS pour un design séquentiel;
  \item routage global/détaillé et optimisation post-route;
  \item extraction RC, timing, connectivité et DRC;
  \item export des résultats qualifiés.
\end{enumerate}

\begin{keybox}
Chaque étape possède son propre gate. Par exemple, placement légal, CTS,
routage, timing, connectivité signal, connectivité spéciale et DRC doivent
rester des résultats distincts.
\end{keybox}

\chapter{Préflight avant Innovus}

\section{Contrôles minimaux}

Avant d'ouvrir Innovus, vérifier :

\begin{itemize}
  \item la présence et la taille non nulle de la netlist et du SDC Genus;
  \item le nom exact du top dans la netlist;
  \item l'absence de modules non résolus;
  \item la présence des macros utilisées dans les LEF;
  \item la compatibilité des noms de cellules entre netlist, Liberty et LEF;
  \item le site de placement et les noms des pins PG des cellules;
  \item les noms des réseaux globaux, par exemple \texttt{VDD} et \texttt{VSS};
  \item la cohérence entre les corners Genus et Innovus.
\end{itemize}

\section{Variables de laboratoire}

Les chemins privés peuvent être fournis par l'environnement :

\begin{lstlisting}[style=guidecode,language=bash]
export TECH_LEF=/chemin/autorise/technology.lef
export STDCELL_LEF=/chemin/autorise/stdcells.lef
export LIB_TC=/chemin/autorise/stdcells_tc.lib
export QRC_TC=/chemin/autorise/qrc_tc.tch
export STDCELL_SITE=core_site
export POWER_NET=VDD
export GROUND_NET=VSS
export STDCELL_POWER_PIN=vddi
export STDCELL_GROUND_PIN=gndi
\end{lstlisting}

Ne pas supposer que \texttt{VDD}, \texttt{VSS}, \texttt{vddi} ou
\texttt{gndi} sont universels. Les vérifier dans le LEF et dans la méthode du
laboratoire.

\chapter{Flow minimal pour le full adder}

\section{Commande directe}

\begin{lstlisting}[style=guidecode,language=bash]
export PNR_NETLIST=/chemin/full_adder_comb.mapped.v
export PNR_SDC=/chemin/full_adder_comb.mapped.sdc
export PNR_RUN_DIR=/tmp/digi_tuto_runs/innovus/full_adder_comb/run_01
innovus -no_gui -files innovus_minimal.tcl
\end{lstlisting}

\section{Script Tcl minimal}

Le script suivant montre import, floorplan, placement, routage et export. Il
utilise une seule vue typique pour rester lisible.

\begin{lstlisting}[style=guidecode,language=TclGuide,
caption={innovus\_minimal.tcl}]
proc require_env {name} {
    if {![info exists ::env($name)] ||
        [string trim $::env($name)] eq ""} {
        error "variable obligatoire absente: $name"
    }
    return $::env($name)
}

set TOP full_adder_comb
set RUN_DIR [file normalize [require_env PNR_RUN_DIR]]
set REPORTS [file join $RUN_DIR reports]
set OUTPUTS [file join $RUN_DIR outputs]
set GENERATED [file join $RUN_DIR generated]
set CHECKPOINTS [file join $RUN_DIR checkpoints]
foreach dir [list $REPORTS $OUTPUTS $GENERATED $CHECKPOINTS] {
    file mkdir $dir
}

# Création locale de la vue TC.
set MMMC [file join $GENERATED basic_tc.mmmc.tcl]
set fh [open $MMMC w]
puts $fh [list create_constraint_mode -name FUNC \
    -sdc_files [list [require_env PNR_SDC]]]
puts $fh [list create_library_set -name LIBSET_TC \
    -timing [list [require_env LIB_TC]]]
puts $fh [list create_rc_corner -name RC_TC \
    -temperature 25 \
    -qx_tech_file [require_env QRC_TC]]
puts $fh [list create_delay_corner -name CORNER_TC \
    -library_set LIBSET_TC -rc_corner RC_TC]
puts $fh [list create_analysis_view -name VIEW_TC \
    -constraint_mode FUNC -delay_corner CORNER_TC]
puts $fh [list set_analysis_view \
    -setup [list VIEW_TC] -hold [list VIEW_TC]]
close $fh

global init_top_cell init_verilog init_lef_file init_mmmc_file
global init_pwr_net init_gnd_net
set init_top_cell $TOP
set init_verilog [require_env PNR_NETLIST]
set init_lef_file [list \
    [require_env TECH_LEF] \
    [require_env STDCELL_LEF]]
set init_mmmc_file $MMMC
set init_pwr_net [require_env POWER_NET]
set init_gnd_net [require_env GROUND_NET]

init_design

globalNetConnect $init_pwr_net -type pgpin \
    -pin [require_env STDCELL_POWER_PIN] -inst * -verbose
globalNetConnect $init_gnd_net -type pgpin \
    -pin [require_env STDCELL_GROUND_PIN] -inst * -verbose
applyGlobalNets

redirect -file [file join $REPORTS check_design_import.rpt] {
    checkDesign -all
}

floorPlan -site [require_env STDCELL_SITE] \
    -r 1.0 0.60 20.0 20.0 20.0 20.0
defOut -floorplan [file join $OUTPUTS ${TOP}.floorplan.def]

placeDesign
redirect -file [file join $REPORTS check_place.rpt] {
    checkPlace
}

# Adapter les couches à la technologie réelle.
setDesignMode -bottomRoutingLayer MET1 -topRoutingLayer MET3
routeDesign -placementCheck
saveDesign [file join $CHECKPOINTS basic_routed.enc]

redirect -file [file join $REPORTS connectivity_regular.rpt] {
    verifyConnectivity -type regular -error 1000 -warning 1000
}
redirect -file [file join $REPORTS drc.rpt] {
    verify_drc
}

set DEF_OUT [file join $OUTPUTS ${TOP}.routed.def]
set VERILOG_OUT [file join $OUTPUTS ${TOP}.routed.v]
defOut $DEF_OUT
saveNetlist $VERILOG_OUT

foreach path [list $DEF_OUT $VERILOG_OUT] {
    if {![file exists $path] || [file size $path] == 0} {
        error "sortie absente ou vide: $path"
    }
}

puts "DIGI_PNR_STATUS: BASIC_FLOW_COMPLETED"
puts "SIGNOFF_STATUS: NOT_ACHIEVED"
exit 0
\end{lstlisting}

\section{Ce que ce flow prouve}

Il prouve que les données peuvent être importées et qu'un premier placement et
routage sont exécutables. Il ne prouve pas encore :

\begin{itemize}
  \item un réseau d'alimentation complet;
  \item le CTS d'un bloc séquentiel;
  \item la fermeture de tous les corners;
  \item zéro violation DRC ou connectivité;
  \item une exportation GDS qualifiée.
\end{itemize}

\chapter{Floorplan et réseau d'alimentation}

\section{Dimensionner le core}

Si l'aire totale des cellules est $A_{cells}$ et l'utilisation visée $U$ :
\[
  A_{core} \geq \frac{A_{cells}}{U}.
\]
Une utilisation de \(0.60\) signifie \(60\,\%\) de cellules et environ
\(40\,\%\) de marge brute pour le routage, les buffers, les fillers et les
optimisations. Cette marge n'est pas une garantie de routabilité.

Le floorplan doit respecter le site, la grille, les marges de core, les
dimensions minimales et l'accès aux pins.

\section{Pins d'interface}

Un placement de pins utile doit :

\begin{itemize}
  \item suivre la connectivité attendue au niveau supérieur;
  \item utiliser une couche et une direction autorisées;
  \item respecter l'espacement et l'accès routable;
  \item garder clock, reset et bus clairement identifiables;
  \item être contrôlé après floorplan et après routage.
\end{itemize}

\section{PG simple}

\begin{lstlisting}[style=guidecode,language=TclGuide]
globalNetConnect VDD -type pgpin \
    -pin vddi -inst * -verbose
globalNetConnect VSS -type pgpin \
    -pin gndi -inst * -verbose
applyGlobalNets

# Valeurs et couches à remplacer par celles du laboratoire.
addRing -nets {VDD VSS} \
    -type core_rings \
    -layer {top MET3 bottom MET3 left MET2 right MET2} \
    -width 2.0 -spacing 1.0 -offset 1.0

sroute -connect {corePin blockPin padPin padRing floatingStripe} \
    -nets {VDD VSS}

redirect -file reports/connectivity_special.rpt {
    verifyConnectivity -type special -error 1000 -warning 1000
}
\end{lstlisting}

\noindent\textbf{\color{GuideOrange}Contrôle PG.}
La présence visuelle d'un ring ou la réussite de \texttt{sroute} ne suffit
pas : lire le rapport de connectivité spéciale et vérifier séparément les nets
d'alimentation et de masse.

\chapter{Placement et optimisation}

\section{Placement}

\begin{lstlisting}[style=guidecode,language=TclGuide]
setPlaceMode -place_global_cong_effort high
placeDesign

redirect -file reports/check_place.rpt {
    checkPlace
}
redirect -file reports/timing_pre_cts.rpt {
    timeDesign -preCTS
}
\end{lstlisting}

Le rapport de placement doit être lu pour les overlaps, cellules hors rows,
sites invalides, orientations et densité locale. Le timing pré-CTS utilise un
modèle de clock idéal; il doit être identifié comme tel.

\section{Optimisation pré-CTS}

\begin{lstlisting}[style=guidecode,language=TclGuide]
optDesign -preCTS
redirect -file reports/timing_pre_cts_optimized.rpt {
    timeDesign -preCTS
}
\end{lstlisting}

Comparer avant/après : WNS, TNS, nombre de violations, aire, congestion et
design rules. Une amélioration du WNS qui crée une congestion sévère n'est pas
une fermeture.

\chapter{CTS pour un design séquentiel}

Le full adder combinatoire n'a pas de clock et saute cette étape. Pour
l'additionneur pipeliné :

\begin{lstlisting}[style=guidecode,language=TclGuide]
set_ccopt_property buffer_cells \
    {CLKBUF_X2 CLKBUF_X4 CLKBUF_X8}
set_ccopt_property inverter_cells \
    {CLKINV_X2 CLKINV_X4}

create_ccopt_clock_tree_spec \
    -file generated/ccopt.spec
source generated/ccopt.spec
ccopt_design

redirect -file reports/clock_tree.rpt {
    report_ccopt_clock_trees
}
redirect -file reports/timing_post_cts.rpt {
    timeDesign -postCTS
}
\end{lstlisting}

Les noms de cellules CTS doivent être validés dans la bibliothèque. Examiner
latence, skew, transitions, sinks non connectés, clock gating et timing setup
et hold. L'optimisation post-CTS s'exécute après ces contrôles :

\begin{lstlisting}[style=guidecode,language=TclGuide]
optDesign -postCTS
optDesign -postCTS -hold
\end{lstlisting}

\chapter{Routage et vérifications}

\section{Couches et routage}

\begin{lstlisting}[style=guidecode,language=TclGuide]
# Limites à adapter au stack métallique du PDK.
setDesignMode -bottomRoutingLayer MET1 -topRoutingLayer MET3

routeDesign
optDesign -postRoute
optDesign -postRoute -hold
saveDesign checkpoints/post_route.enc
\end{lstlisting}

Une limite de couche trop basse peut créer congestion et détours; une limite
trop haute peut violer le contrat d'intégration du bloc. Le choix doit être
écrit dans la méthodologie de handoff.

\section{Contrôles physiques séparés}

\begin{lstlisting}[style=guidecode,language=TclGuide]
redirect -file reports/connectivity_regular.rpt {
    verifyConnectivity -type regular -error 1000 -warning 1000
}
redirect -file reports/connectivity_special.rpt {
    verifyConnectivity -type special -error 1000 -warning 1000
}
redirect -file reports/drc.rpt {
    verify_drc
}
redirect -file reports/antenna.rpt {
    verifyProcessAntenna
}
\end{lstlisting}

\begin{longtable}{p{0.26\textwidth}p{0.64\textwidth}}
\toprule
Gate & Question \\
\midrule
Connectivité regular & Tous les signaux ont-ils un chemin électrique complet,
sans open ni short? \\
Connectivité special & VDD et VSS atteignent-ils correctement rings, rails et
pins PG? \\
DRC & Les géométries respectent-elles les règles vérifiées par Innovus? \\
Antenne & Les ratios d'antenne sont-ils acceptables ou réparés? \\
Timing & Setup et hold sont-ils fermés dans chaque vue requise? \\
\bottomrule
\end{longtable}

\section{Extraction et timing post-route}

\begin{lstlisting}[style=guidecode,language=TclGuide]
extractRC

redirect -file reports/timing_post_route_setup.rpt {
    timeDesign -postRoute
}
redirect -file reports/timing_post_route_hold.rpt {
    timeDesign -postRoute -hold
}

rcOut -spef outputs/design.post_route.spef
\end{lstlisting}

Le timing post-route doit utiliser les vues et fichiers QRC attendus. Vérifier
les clocks propagées, les chemins setup/hold, les violations électriques, les
chemins non contraints et la cohérence des unités.

\chapter{Wrapper de lancement}

Ce wrapper effectue un préflight simple, crée un dossier unique et garde le
code retour Innovus.

\begin{lstlisting}[style=guidecode,language=bash,
caption={run\_innovus.sh}]
#!/usr/bin/env bash
set -Eeuo pipefail

INNOVUS_BIN=${INNOVUS_BIN:-innovus}
RUN_ROOT=${RUN_ROOT:-/tmp/digi_tuto_runs}
RUN_ID=${RUN_ID:-$(date -u +%Y%m%dT%H%M%SZ)}
export PNR_RUN_DIR="$RUN_ROOT/innovus/full_adder_comb/$RUN_ID"

required_vars=(
  PNR_NETLIST PNR_SDC TECH_LEF STDCELL_LEF
  LIB_TC QRC_TC STDCELL_SITE POWER_NET GROUND_NET
  STDCELL_POWER_PIN STDCELL_GROUND_PIN
)
for name in "${required_vars[@]}"; do
  [[ -n ${!name:-} ]] || {
    echo "ERREUR: variable absente: $name" >&2
    exit 2
  }
done

for path in \
  "$PNR_NETLIST" "$PNR_SDC" "$TECH_LEF" \
  "$STDCELL_LEF" "$LIB_TC" "$QRC_TC" \
  innovus_minimal.tcl; do
  [[ -s "$path" ]] || {
    echo "ERREUR: fichier absent ou vide: $path" >&2
    exit 2
  }
done

command -v "$INNOVUS_BIN" >/dev/null 2>&1 || {
  echo "ERREUR: Innovus introuvable" >&2
  exit 127
}

[[ ! -e "$PNR_RUN_DIR" ]] || {
  echo "ERREUR: dossier déjà présent: $PNR_RUN_DIR" >&2
  exit 2
}
mkdir -p "$PNR_RUN_DIR/logs"

set +e
"$INNOVUS_BIN" -no_gui -files innovus_minimal.tcl \
  -log "$PNR_RUN_DIR/logs/innovus.log"
RC=$?
set -e

(( RC == 0 )) || {
  echo "RESULTAT: FAIL (innovus=$RC)" >&2
  exit "$RC"
}

for path in \
  "$PNR_RUN_DIR/outputs/full_adder_comb.routed.def" \
  "$PNR_RUN_DIR/outputs/full_adder_comb.routed.v" \
  "$PNR_RUN_DIR/reports/drc.rpt" \
  "$PNR_RUN_DIR/reports/connectivity_regular.rpt"; do
  [[ -s "$path" ]] || {
    echo "RESULTAT: FAIL (sortie absente: $path)" >&2
    exit 3
  }
done

echo "RESULTAT: FLOW_EXECUTE"
echo "Lire tous les rapports avant de qualifier le résultat"
\end{lstlisting}

\chapter{Méthodologie avancée MMMC}

\section{Définition des vues}

\begin{lstlisting}[style=guidecode,language=TclGuide]
create_constraint_mode -name FUNC -sdc_files [list $PNR_SDC]

create_library_set -name LIB_BC -timing [list $LIBERTY_BC]
create_library_set -name LIB_TC -timing [list $LIBERTY_TC]
create_library_set -name LIB_WC -timing [list $LIBERTY_WC]

create_rc_corner -name RC_BC -qx_tech_file $QRC_BC
create_rc_corner -name RC_TC -qx_tech_file $QRC_TC
create_rc_corner -name RC_WC -qx_tech_file $QRC_WC

create_delay_corner -name DC_BC \
    -library_set LIB_BC -rc_corner RC_BC
create_delay_corner -name DC_TC \
    -library_set LIB_TC -rc_corner RC_TC
create_delay_corner -name DC_WC \
    -library_set LIB_WC -rc_corner RC_WC

create_analysis_view -name VIEW_BC \
    -constraint_mode FUNC -delay_corner DC_BC
create_analysis_view -name VIEW_TC \
    -constraint_mode FUNC -delay_corner DC_TC
create_analysis_view -name VIEW_WC \
    -constraint_mode FUNC -delay_corner DC_WC

set_analysis_view \
    -setup [list VIEW_TC VIEW_WC] \
    -hold  [list VIEW_TC VIEW_BC]
\end{lstlisting}

Le PDK doit définir températures, tension, modèles RC et association correcte
des vues. Vérifier dans Innovus les vues actives avant chaque rapport.

\section{Organisation en stages}

Un flow réutilisable sépare les fichiers Tcl :

\begin{enumerate}
  \item import et validation de la technologie;
  \item floorplan, pins et PG;
  \item placement, optimisation et CTS conditionnel;
  \item routage, extraction, vérifications et exports.
\end{enumerate}

Chaque stage doit :

\begin{itemize}
  \item retourner une erreur explicite en cas d'échec;
  \item produire ses propres rapports;
  \item sauvegarder un checkpoint utile;
  \item ne marquer PASS qu'après vérification de ses sorties;
  \item conserver les gates non atteints comme tels.
\end{itemize}

\chapter{Exports et niveau de confiance}

\section{Exports}

\begin{lstlisting}[style=guidecode,language=TclGuide]
defOut outputs/design.routed.def
saveNetlist outputs/design.routed.v
write_sdc -view VIEW_WC > outputs/design.routed.sdc
rcOut -spef outputs/design.post_route.spef
saveDesign checkpoints/final.enc

# Uniquement lorsque stream map et GDS de référence sont validés.
streamOut outputs/design.gds \
    -mapFile $STREAM_MAP \
    -merge $STDCELL_GDS \
    -units 1000 \
    -mode ALL
\end{lstlisting}

Un GDS créé est un candidat d'export, pas une preuve de signoff. Les checks
foundry et la comparaison layout-vers-source appartiennent à une qualification
séparée avec les outils et decks autorisés.

\section{Matrice de statut}

\begin{longtable}{p{0.26\textwidth}p{0.23\textwidth}p{0.41\textwidth}}
\toprule
Gate & Valeur & Preuve attendue \\
\midrule
Préflight & PASS/FAIL & fichiers et variables contrôlés \\
Import & PASS/FAIL & \texttt{checkDesign}, cellules et vues cohérentes \\
Floorplan & PASS/FAIL & dimensions, rows, pins et PG examinés \\
Placement & PASS/FAIL & \texttt{checkPlace} sans erreur bloquante \\
CTS & PASS/FAIL/NON APPLICABLE & arbre et sinks qualifiés \\
Routage & PASS/FAIL & routage terminé et checkpoint lisible \\
Connectivité signal & PASS/FAIL & rapport regular explicite \\
Connectivité PG & PASS/FAIL & rapport special explicite \\
DRC Innovus & PASS/FAIL & total de violations connu \\
Antenne & PASS/FAIL/NON EXÉCUTÉ & rapport et total connus \\
Timing & PASS/FAIL & setup et hold par vue \\
Export & PASS/FAIL & fichiers attendus non vides \\
Signoff & NON ATTEINT/ATTEINT & contrôles externes requis terminés \\
\bottomrule
\end{longtable}

\chapter{Diagnostic}

\begin{longtable}{p{0.22\textwidth}p{0.31\textwidth}p{0.37\textwidth}}
\toprule
Symptôme & Cause probable & Première action \\
\midrule
Cellule introuvable & désaccord netlist/LEF/Liberty &
Comparer le nom exact dans les trois entrées. \\
Site invalide & mauvais \texttt{STDCELL\_SITE} &
Lire le nom du site dans le LEF. \\
Pins PG non connectées & mauvais noms de pins ou nets &
Contrôler LEF, \texttt{globalNetConnect} et rapport special. \\
Placement illégal & rows, orientation ou dimensions &
Lire la première erreur de \texttt{checkPlace}. \\
CTS sans sinks & clock SDC absente ou mauvais objet &
Vérifier la clock et les registres après import. \\
Routage impossible & congestion ou limites de couches &
Lire la congestion et le premier type de violation. \\
Opens signal & accès pin ou route incomplète &
Localiser le premier net du rapport regular. \\
Shorts PG & ring, stripe ou via incompatible &
Analyser le premier marker et la connectivité special. \\
Timing post-route négatif & RC ou chemin critique réel &
Lire le chemin complet et la vue avant optimisation. \\
GDS vide ou incomplet & map ou merge incorrect &
Vérifier fichiers d'entrée, taille et messages de streamout. \\
\bottomrule
\end{longtable}

\begin{failurebox}[Ordre de debug]
Corriger d'abord l'import, puis le floorplan/PG, le placement, le CTS, le
routage, la connectivité et enfin le timing post-route. Une optimisation ne
doit pas masquer un problème d'entrée ou de topologie.
\end{failurebox}

\appendix
\chapter{Checklist finale}

\begin{enumerate}
  \item Netlist, SDC, LEF, Liberty et QRC sont présents et cohérents.
  \item Le top, le site, les nets PG et les pins PG sont vérifiés.
  \item Les vues setup/hold actives correspondent à la méthode du PDK.
  \item Floorplan, rows et pins sont contrôlés.
  \item La connectivité spéciale est vérifiée après création du PG.
  \item Le placement est légal et la congestion examinée.
  \item Le CTS est qualifié ou explicitement non applicable.
  \item Routage, connectivité signal, connectivité PG, DRC et antenne ont des
    statuts séparés.
  \item Le timing post-route setup/hold est lu dans chaque vue requise.
  \item DEF, netlist, SDC, SPEF, checkpoint et éventuel GDS sont non vides.
  \item Aucun export n'est appelé « signoff » sans les preuves correspondantes.
\end{enumerate}

\end{document}
